\section*{Chapitre \ref{ch:qm:floating-point-and-numerical-linear-algebra} --- Virgule flottante et algèbre linéaire numérique}

\begin{solution}{exo:qm:floating-point-and-numerical-linear-algebra:1}
Unité d'arrondi $u = 2^{-53} \approx 1.11 \times 10^{-16}$, \omterm{def:qm:floating-point-and-numerical-linear-algebra:fp}{epsilon machine} $2^{-52} \approx 2.22 \times 10^{-16}$~; $53\log_{10}2 \approx 15.95$, soit quinze à seize chiffres décimaux significatifs.
\end{solution}

\begin{solution}{exo:qm:floating-point-and-numerical-linear-algebra:2}
Un demi-centième par recalcul~: $0.005 \times 10\,000 \times 250 = 12\,500$ points par an, à la baisse, avant la capitalisation avec les propres variations de l'indice.
\end{solution}

\begin{solution}{exo:qm:floating-point-and-numerical-linear-algebra:3}
$10^8 + 1$ demande 27 bits et est représentable, donc les deux opérations sont exactes. Au-delà de $2^{53} \approx 9.0 \times 10^{15}$, l'écart entre deux nombres en double précision consécutifs vaut 2, donc $10^{16} + 1$ est un cas d'égalité, arrondi au
voisin dont la mantisse est paire, qui est $10^{16}$~: le résultat vaut 0.
\end{solution}

\begin{solution}{exo:qm:floating-point-and-numerical-linear-algebra:4}
Chacun de $\sum x_i^2$ et $(\sum x_i)^2/n$ vaut environ $nc^2$ et porte une erreur d'arrondi d'au moins $u\,nc^2$ (davantage avec une sommation récursive)~; leur différence, divisée par $n - 1$, a une erreur absolue
de l'ordre de $uc^2$. Avec $c = 10^4$~: $1.1 \times 10^{-16} \times 10^8 \approx 10^{-8}$, soit environ $10^{-7}$ relativement à une variance de 0.08~; l'erreur de sommation la fait remonter, et le tutoriel
mesure $8.0 \times 10^{-7}$.
\end{solution}

\begin{solution}{exo:qm:floating-point-and-numerical-linear-algebra:5}
Les équations normales travaillent avec $\kappa^2 = 10^{10}$ et peuvent perdre environ dix des seize chiffres~; la \omterm{def:qm:floating-point-and-numerical-linear-algebra:factor}{factorisation QR}, avec $\kappa = 10^5$, en perd environ cinq.
\end{solution}

\begin{solution}{exo:qm:floating-point-and-numerical-linear-algebra:6}
Tous les $x_i - m$ sont au plus égaux à 0, donc chaque exponentielle est au plus égale à 1 et la plus grande vaut exactement 1~: la somme appartient à $[1, n]$ et son logarithme ne peut ni déborder ni valoir $-\infty$. $\log(e^{1000} + e^{1000}) =
1000 + \ln2 = 1000.693$, là où la formule directe déborde vers l'infini.
\end{solution}

\begin{solution}{exo:qm:floating-point-and-numerical-linear-algebra:7}
\texttt{summation\_errors()}~: pour $n = 10^6$, les erreurs par rapport à la somme rationnelle exacte (47\,878\,411.409) valent $2.5 \times 10^{-6}$ (naïve), $2.8 \times 10^{-8}$ (par paires) et $2.2 \times 10^{-9}$
(Neumaier), la dernière étant inférieure à une \omterm{def:qm:floating-point-and-numerical-linear-algebra:fp}{unité de dernière position} (ulp, soit $7.5 \times 10^{-9}$).
\end{solution}

\begin{solution}{exo:qm:floating-point-and-numerical-linear-algebra:8}
Un écart au dixième chiffre est ce que produisent des ordres d'opérations différents, une contraction par \omterm{def:qm:floating-point-and-numerical-linear-algebra:fma}{multiplication-addition fusionnée} (FMA) ou des implémentations différentes de \texttt{exp} et \texttt{log} sur un calcul modérément
conditionné~; ce n'est pas la preuve d'un bogue. Une preuve serait un écart supérieur au \omterm{def:qm:floating-point-and-numerical-linear-algebra:cond}{conditionnement} multiplié par l'unité d'arrondi, ou un écart sur des données de référence où les deux moteurs sont censés
effectuer des opérations identiques. Soit on fixe l'ordre et les options et l'on vérifie les bits, soit on compare avec une tolérance justifiée par le \omterm{def:qm:floating-point-and-numerical-linear-algebra:cond}{conditionnement}.
\end{solution}

\begin{solution}{pb:qm:floating-point-and-numerical-linear-algebra:1}
\textbf{1.} Si les chiffres au-delà de la troisième décimale sont répartis uniformément, la partie écartée est uniforme sur $[0, 0.001)$~: 0.0005 en moyenne, toujours une perte.
\textbf{2.} $0.0005 \times 2\,400 \times 480 = 576$ points, à comparer avec $1098.892 - 524.811 = 574.081$.
\textbf{3.} $0.0005 \times 3\,000 \times 480 = 720$ points.
\textbf{4.} Les erreurs d'arrondi sont symétriques autour de zéro~: pas de dérive, seulement un bruit d'écart-type voisin de $0.001\sqrt{n/12} \approx 0.3$~; l'indice arrondi simulé termine à 1098.57, contre 1098.892.
\textbf{5.} À 422.23, une perte de 676.7. Un point perdu à l'étape $k$ est reporté par la croissance de l'indice entre $k$ et la fin, si bien que la perte vaut $0.0005\sum_kE_n/E_k$~; l'indice exact simulé
passe l'essentiel de son temps sous sa valeur finale (entre 803 et 1183), et le rapport moyen vaut 1.175, donc $576 \times 1.175 = 676.7$.
\textbf{6.} $2.5 \times 10^{-6}$, $2.8 \times 10^{-8}$ et $2.2 \times 10^{-9}$.
\textbf{7.} Erreurs relatives $8.0 \times 10^{-7}$, 1.3~\% et 1060 (86.84 pour une valeur exacte de 0.0819).
\textbf{8.} Elles retranchent la moyenne avant d'élever au carré, donc elles élèvent au carré de petits écarts~; leurs erreurs sont de l'ordre de $u$ fois la dispersion, et non de $u$ fois le carré du décalage (environ $10^{-12}$ au plus pour deux
passages, $1.6 \times 10^{-8}$ pour Welford avec un décalage de $10^8$).
\textbf{9.} Erreurs relatives $5.3 \times 10^{-3}$ et de l'ordre de $10^{-10}$ (ses chiffres exacts relèvent de l'arrondi et varient selon la machine)~: les équations normales conservent environ deux chiffres, la \omterm{def:qm:floating-point-and-numerical-linear-algebra:factor}{factorisation QR} environ dix (sept perdus, comme le prédit $\kappa = 10^7$).
\textbf{10.} Au sixième pivot (indice 5), qui vaut $-0.14$~: la matrice estimée par paires n'est pas semi-définie positive (plus petite valeur propre $-1.42$), donc d'après
la \cref{prop:qm:floating-point-and-numerical-linear-algebra:chol}, la factorisation doit échouer.
\textbf{11.} Les bits de huit résultats (les sommes naïve, par paires et de Neumaier, la moyenne et la variance de Welford, un log-sum-exp, les première et dernière composantes d'une résolution de système tridiagonal), vérifiés comme
les mêmes motifs hexadécimaux dans les tests Python, C++20 et Rust.
\textbf{12.} Deux arrondis donnent exactement 0~; un seul arrondi donne $2^{-54} \approx 5.6 \times 10^{-17}$, l'erreur de représentation de $0.1$ multipliée par 10.
\textbf{13.} \texttt{-ffp-contract=fast}, valeur par défaut de GCC hors d'un mode C conforme aux normes, dès que la cible (\texttt{-march}) dispose d'une \omterm{def:qm:floating-point-and-numerical-linear-algebra:fma}{multiplication-addition fusionnée}~; aussi \texttt{-ffast-math}.
\textbf{14.} Environ 1\,900 sans \omterm{def:qm:floating-point-and-numerical-linear-algebra:cg}{préconditionneur} (1\,884 avec un seul fil d'exécution BLAS~; le nombre dépend de l'ordre de sommation) et 102 avec le \omterm{def:qm:floating-point-and-numerical-linear-algebra:cg}{préconditionneur}.
\textbf{15.} L'arrondi détruit la conjugaison des directions, si bien que l'argument de terminaison finie tombe~; ce qui subsiste est la vitesse en $\sqrt\kappa$ ($\sqrt{1.1 \times 10^5} \approx 330$), dont la borne
prévoit environ 3\,100 itérations pour un facteur $10^8$~; les quelque 1\,900 observées s'y inscrivent.
\textbf{16.} Recalculer entièrement l'indice, à partir des cours des composantes et du diviseur, une fois par jour, et le comparer à la valeur courante~: un écart croissant d'environ un point par jour est visible en une semaine.
\textbf{17.} Empreintes de hachage des données d'entrée, versions du code et des bibliothèques, compilateur et options, nombre de fils d'exécution et ordre de réduction, graines aléatoires, et motifs binaires des agrégats intermédiaires clés.
\textbf{18.} Lorsque des chiffres sont réconciliés entre systèmes ou dans le temps (limites, états réglementaires, tests de non-régression d'un backtest), là où un écart inexpliqué coûte du temps d'investigation~; le coût tient
aux réductions à ordre fixe et à l'absence de fast-math.
\textbf{19.} Résultat nommé~: \emph{l'indice qui a perdu la moitié de sa valeur}~: la troncature retire 0.0005 par recalcul en moyenne, soit $0.0005 \times 2\,400 \times 480 = 576$ points, à comparer avec un manque documenté
de 574.081 (524.811 contre 1098.892).
\textbf{20.} La machine l'arrondit à chaque opération, avec une erreur relative d'au plus $2^{-53}$, et ce que l'algorithme et le \omterm{def:qm:floating-point-and-numerical-linear-algebra:cond}{conditionnement} font de ces arrondis décide de la réponse.
\end{solution}

\begin{solution}{iq:qm:floating-point-and-numerical-linear-algebra:1}
Ni 0.1, ni 0.2, ni 0.3 ne sont représentables en binaire~; la somme arrondie des deux premiers dépasse d'une \omterm{def:qm:floating-point-and-numerical-linear-algebra:fp}{unité de dernière position} (ulp) le 0.3 arrondi. Comparer avec une tolérance mêlant une partie relative et une partie absolue,
$|a - b| \le \max(r\max(|a|, |b|), t)$, choisie selon le \omterm{def:qm:floating-point-and-numerical-linear-algebra:cond}{conditionnement} du calcul~; ou travailler en entiers (ticks, centimes) lorsque l'exactitude est requise.
\iqlookfor{Représentation binaire, et tolérance mêlant partie relative et partie absolue, ou unités entières.}
\end{solution}

\begin{solution}{iq:qm:floating-point-and-numerical-linear-algebra:2}
Des ordres d'opérations différents (réductions vectorisées ou parallèles, nombres de fils d'exécution), contraction par \omterm{def:qm:floating-point-and-numerical-linear-algebra:fma}{multiplication-addition fusionnée}, options de fast-math, implémentations différentes des fonctions transcendantes (\texttt{exp}, \texttt{log},
\texttt{pow} ne sont pas arrondies correctement), précision étendue dans les registres intermédiaires, analyse différente des entrées décimales, et algorithmes de bibliothèque différents (un BLAS face à une boucle écrite à la main).
\iqlookfor{Non-associativité, et au moins trois sources concrètes.}
\end{solution}

\begin{solution}{iq:qm:floating-point-and-numerical-linear-algebra:3}
Mise à jour de Welford~: $\delta = x - m$, $m \leftarrow m + \delta/n$, $M_2 \leftarrow M_2 + \delta(x - m)$. La différence entre la somme des carrés et le carré de la somme subit une \omterm{def:qm:floating-point-and-numerical-linear-algebra:cancel}{annulation catastrophique} lorsque la moyenne est grande
par rapport à la dispersion, comme c'est le cas des prix~; pour des prix voisins de $10^8$, elle donne des résultats absurdes.
\iqlookfor{Récurrence de Welford et argument d'annulation.}
\end{solution}

\begin{solution}{iq:qm:floating-point-and-numerical-linear-algebra:4}
Les équations normales forment $X^\top X$, dont le \omterm{def:qm:floating-point-and-numerical-linear-algebra:cond}{conditionnement} vaut $\kappa(X)^2$, et perdent deux fois plus de chiffres~; la \omterm{def:qm:floating-point-and-numerical-linear-algebra:factor}{factorisation QR} travaille directement sur $X$ et est à \omterm{def:qm:floating-point-and-numerical-linear-algebra:cond}{stabilité inverse}. Pour $\kappa = 10^7$, les coefficients du chapitre
sont faux dès le troisième chiffre avec les équations normales et justes à dix chiffres avec la QR.
\iqlookfor{Carré du \omterm{def:qm:floating-point-and-numerical-linear-algebra:cond}{conditionnement}.}
\end{solution}

\begin{solution}{iq:qm:floating-point-and-numerical-linear-algebra:5}
Commencer par se demander pourquoi~: une matrice estimée par paires ou avec des données manquantes peut être indéfinie, et une \omterm{def:qm:covariance-estimation-and-random-matrices:cov}{matrice de covariance empirique} comportant moins d'observations que d'actifs est singulière. Regarder l'indice défaillant et les
valeurs propres. Puis réparer avec la \omterm{def:qm:convex-optimisation:ncm}{matrice de corrélation la plus proche}, un shrinkage ou un \omterm{def:qm:covariance-estimation-and-random-matrices:factor}{modèle factoriel}~; une petite perturbation (jitter) n'est acceptable que pour une matrice semi-définie positive à l'arrondi près. Ne jamais
la perturber en silence.
\iqlookfor{Diagnostic avant réparation, et réparation fondée sur des principes.}
\end{solution}

\begin{solution}{iq:qm:floating-point-and-numerical-linear-algebra:6}
Fixer le découpage des données en blocs indépendamment du nombre de fils d'exécution, sommer chaque bloc dans un ordre fixe (avec \omterm{def:qm:floating-point-and-numerical-linear-algebra:cancel}{sommation compensée} pour la précision), et combiner les sommes partielles selon un arbre
à ordre fixe~; les fils d'exécution se contentent alors de planifier les blocs. Alternatives~: accumulation exacte (accumulateur long ou virgule fixe entière), ou schéma de pré-arrondi qui rend l'addition exacte.
\iqlookfor{Ordre indépendant du nombre de fils d'exécution.}
\end{solution}
